\chapter{Программная реализация алгоритма маршрутизации}\label{appendix-implementation}
\addcontentsline{toc}{chapter}{Приложение А. Программная реализация алгоритма маршрутизации}

Данное приложение содержит практические аспекты реализации разработанного алгоритма маршрутизации для самоорганизующихся сетей с низкой пропускной способностью.

\section{Требования к программному обеспечению}

Для работы с реализацией алгоритма необходимо:

\begin{itemize}
    \item Операционная система: Linux (рекомендуется Ubuntu 20.04+) или Windows 10+
    \item Установленные пакеты:
    \begin{itemize}
        \item Python 3.8+ (с библиотеками NumPy, SciPy)
        \item OMNeT++ 6.0+ с модулем INET
        \item Git для управления версиями
    \end{itemize}
\end{itemize}

\section{Инструкция по установке}

\begin{enumerate}
    \item Клонировать репозиторий проекта:
    \begin{verbatim}
    git clone https://github.com/example/ad-hoc-routing.git
    \end{verbatim}
    
    \item Установить зависимости:
    \begin{verbatim}
    pip install -r requirements.txt
    \end{verbatim}
    
    \item Импортировать проект в OMNeT++ IDE:
    \begin{itemize}
        \item File → Import → Existing Projects into Workspace
        \item Выбрать папку проекта
    \end{itemize}
\end{enumerate}

\section{Конфигурация сети}

Основные параметры конфигурации:

\begin{table}[h]
\centering
\begin{tabular}{|l|c|}
\hline
Параметр & Значение \\ \hline
Количество абонентов & 15-50 \\
Скорость передачи & 32-64 бит/с \\
Максимальная скорость абонентов & 60 км/ч \\
Мощность передачи & 10-100 мВт \\ \hline
\end{tabular}
\caption{Параметры конфигурации сети}
\label{tab:network-config}
\end{table}

\section{Примеры реализации ключевых функций}

\subsection{Алгоритм маршрутизации}

\begin{verbatim}
def find_optimal_route(network_graph, current_node, destination):
    # Реализация алгоритма выбора маршрута
    routes = []
    for neighbor in network_graph.neighbors(current_node):
        metric = calculate_route_metric(neighbor)
        routes.append((neighbor, metric))
    
    return min(routes, key=lambda x: x[1])
\end{verbatim}

\subsection{Модуль энергоменеджмента}

\begin{verbatim}
def adjust_power_level(distance, remaining_energy):
    # Адаптивная регулировка мощности
    base_power = 10  # мВт
    max_power = 100  # мВт
    power = base_power + (max_power - base_power) * (distance / max_distance)
    
    if remaining_energy < energy_threshold:
        power *= 0.7  # Снижение мощности при низком заряде
        
    return min(power, max_power)
\end{verbatim}

\section{Рекомендации по использованию}

\begin{itemize}
    \item Для тестирования использовать сценарии из папки \texttt{simulations/}
    \item Параметры сети изменять в файле \texttt{config/network.ini}
    \item Результаты сохраняются в \texttt{results/logs/} в формате CSV
\end{itemize}

% \begin{figure}[h]
% \centering
% \includegraphics[width=0.8\textwidth]{network-scheme}
% \caption{Схема тестовой сети в OMNeT++}
% \label{fig:network-scheme}
% \end{figure}